\documentclass[10pt,a4paper]{amsart}
\usepackage[margin=19mm]{geometry}
\usepackage{amsmath,amssymb,mathtools}
\usepackage[scr=rsfs,cal=euler]{mathalfa}
\usepackage{xfrac}
\usepackage[unicode,hidelinks,bookmarks=false]{hyperref}
\newcommand{\ie}{i.e.\ }
\newcommand{\cf}{cf.\ }
\newcommand{\resp}{respectively\ }
\newcommand{\Subsection}{\subsection}
\providecommand{\nobreakdash}{\nobreak\hbox{-}}
\setlength{\emergencystretch}{2em}
\multlinegap0pt
\usepackage{amssymb}
\usepackage{mathtools}
\usepackage[shortlabels]{enumitem}
\newtheorem{theorem}{Theorem}
\title{\textbf{Schur Visibility and Anti-Phantom Reduction\\ in One-Component Navier--Stokes Degeneration}}
\author{Runlong Yu}
\date{Source: arXiv:2606.12267v1 (CC BY 4.0)}
\begin{document}
\maketitle

\noindent\emph{Source and licence.} \textbf{Schur Visibility and Anti-Phantom Reduction\\ in One-Component Navier--Stokes Degeneration}. Version arXiv:2606.12267v1 (CC BY 4.0), distributed by arXiv under the Creative Commons Attribution 4.0 International licence, http://creativecommons.org/licenses/by/4.0/, as recorded in the arXiv licence metadata for this exact version and shown on the abstract page https://arxiv.org/abs/2606.12267v1. Full licence notice: \texttt{licenses/arxiv-2606.12267-CC-BY-4.0.txt}.\par\smallskip

\noindent\textit{Editorial scope.} Counted unit: the article's theorem eq:tag-10-253 (source0.tex lines 7417--7464). The article's complete original proof of it is delivered with it (source0.tex lines 7466--7534, one \texttt{proof} environment of 69 lines). No mathematics is written, repaired or re-derived: the statement and the proof are the article's own lines, in the article's own notation, and no retained line is substituted or re-flowed.  Retained uncounted as the source-local context the proof needs: lines 7367--7373; lines 7375--7380; lines 7382--7389.  Nothing was omitted from any retained span, so the package carries no omission record.  No retained line contains a citation, so the wrapper carries no bibliography and no bibliography entry.  No retained line contains a figure environment, an image-inclusion command or drawing code, so no image is delivered and the package declares no asset dependency.  Editorial formatting only, in addition: an AMS \texttt{amsart} preamble that restates the article's own declarations and macros used by the retained lines, namely the environments \texttt{theorem} on separate counters, together with the standard packages those lines need.  The retained environments print in this document as Theorem 1 in order of appearance, whereas the article prints the same items with the numbering of its own full document.  The complete native source of the article as distributed in the arXiv e-print for this exact version is bundled unchanged as \texttt{sources/202612/source0.tex}.
\begin{equation}
\langle g_n,y\rangle =
\mathcal T_n(y)
+
\mathcal R_n(y),
\label{eq:tag-10-247}
\end{equation}

\begin{equation}
|\mathcal T_n(y)|
\le
C r_n|A_n^*y|_{H_n},
\label{eq:tag-10-248}
\end{equation}

\begin{equation}
|\mathcal R_n(y)|
\le
C r_n|J_n^{\mathrm{rel}}P_n y|_{W_n}
+
o(r_n)|y|_{Y_n}.
\label{eq:tag-10-249}
\end{equation}

\begin{theorem}[No residual left-singular true phantom in the controlled regime]
Under controlled combined observability and the NS residual representation, no residual left-singular true phantom cascade exists.

Moreover, every NS-derived surviving residual satisfies the relaxed anti-phantom estimate
\begin{equation}
|\langle g_n,y\rangle|
\le
C r_n
\Big(
|A_n^*y|_{H_n}
+
|J_n^{\mathrm{rel}}P_n y|_{W_n}
\Big)
+
o(r_n)|y|_{Y_n}
\label{eq:tag-10-253}
\end{equation}
for all \(y\in Y_n\). In particular, if the relaxed component is kept in the comparison class and controlled by the small vertical component \(u_3\), then the residual has no uncontrolled true phantom projection.

If, in addition, the relaxed component vanishes after cleaning,
\begin{equation}
J_n^{\mathrm{rel}}P_n y=0
\quad\Longrightarrow\quad
P_ny=0,
\label{eq:tag-10-254}
\end{equation}
then the strict vertical-duality estimate follows:
\begin{equation}
|\langle g_n,y\rangle|
\le
C r_n|A_n^*y|_{H_n}
+
o(r_n)|y|_{Y_n}.
\label{eq:tag-10-255}
\end{equation}
Consequently,
\[
g_n\in \operatorname{Range}A_n
\]
up to the \(o(r_n)\) error, and
\begin{equation}
\operatorname{Cost}^{\operatorname{tr}}_{A_n}(g_n)
\le
C r_n^2
\label{eq:tag-10-256}
\end{equation}
after absorbing the error at the branch-native scale.
\end{theorem}

\begin{proof}
The relaxed anti-phantom estimate \eqref{eq:tag-10-253} is immediate from the residual representation. Indeed, combining \eqref{eq:tag-10-247}, \eqref{eq:tag-10-248}, and \eqref{eq:tag-10-249} gives
\[
|\langle g_n,y\rangle|
\le
C r_n|A_n^*y|_{H_n}
+
C r_n|J_n^{\mathrm{rel}}P_n y|_{W_n}
+
o(r_n)|y|_{Y_n}.
\]
This proves \eqref{eq:tag-10-253}.

We now prove that no residual left-singular true phantom cascade exists. Suppose, for contradiction, that such a sequence \(y_n\) exists. Then
\[
|y_n|_{Y_n}=1,
\qquad
M_n\Big(|A_n^*y_n|_{H_n}+|J_n^{\mathrm{rel}}P_n y_n|_{W_n}\Big)\to0.
\]
By the combined observability estimate,
\begin{equation}
1 =
|y_n|_{Y_n}
\le
M_n\Big(
|A_n^*y_n|_{H_n}
+
|J_n^{\mathrm{rel}}P_n y_n|_{W_n}
\Big).
\label{eq:tag-10-257}
\end{equation}
The right-hand side tends to \(0\) by the weighted invisibility condition in the definition of a controlled true phantom cascade. This contradicts \(|y_n|_{Y_n}=1\). Therefore no such cascade exists.

Equivalently, the only way a normalized dual direction can make both
\[
A_n^*y_n
\]
and
\[
J_n^{\mathrm{rel}}P_ny_n
\]
small is by leaving the controlled regime: gauge cleaning fails, localization becomes nonperturbative, homogeneous-tail observability collapses, or the combined observability constant grows faster than the admissible logarithmic scale. These cases have already been isolated as separate failure mechanisms.

Finally assume the stronger cleaned strict condition \eqref{eq:tag-10-254}. If
\[
J_n^{\mathrm{rel}}P_ny=0,
\]
then \(P_ny=0\), so the cleaned phantom component vanishes. Therefore every surviving dual direction is trace-visible, modulo the \(o(r_n)\) controlled errors. From \eqref{eq:tag-10-253} we obtain
\[
|\langle g_n,y\rangle|
\le
C r_n|A_n^*y|_{H_n}
+
o(r_n)|y|_{Y_n}.
\]
After absorbing the \(o(r_n)\)-term into the branch-native scale, this gives
\[
|\langle g_n,y\rangle|
\le
C r_n|A_n^*y|_{H_n}.
\]
By finite-dimensional trace-cost duality,
\[
\operatorname{Cost}^{\operatorname{tr}}_{A_n}(g_n)
\le
C r_n^2.
\]
Equivalently, \(g_n\) lies in \(\operatorname{Range}A_n\) with controlled minimal trace preimage. This is precisely the strict vertical-duality estimate.
\end{proof}

\end{document}
